% GENERATED FILE. Edit canonical lessons, then run npm run generate:books.
% canonical-source-hash: fnv1a64:14746ebe303d6d29
% canonical-lessons: ES-C619-a-lo-mejor, ES-C619-previsible, ES-C619-oportunidad, ES-C619-perspectiva, ES-C619-predecir

\chapter{A lo mejor}
\label{ch:es-a-lo-mejor}
\hlchaptermodality{619}

\begin{chapteropening}
\textbf{By the end of this chapter:} \emph{I can say that something may happen and look ahead.}

Complete the last lesson of chapter 619: I can say that something may happen and look ahead.
\end{chapteropening}

\section[a lo mejor]{a lo mejor — maybe, perhaps}

\label{lesson:ES-C619-a-lo-mejor}



\begin{quote}
Before the new one: say the Spanish for possibly, then the Spanish for most likely.
\end{quote}

\subsection*{You'll want to know: a lo mejor}
\textbf{a lo mejor} — \textquotedblleft{}maybe, perhaps\textquotedblright{}.

\begin{quote}
A lo mejor vamos al cine.
\end{quote}

\begin{quote}
\emph{Maybe we will go to the cinema.}
\end{quote}

Followed by the indicative, while \emph{quizás} often takes the subjunctive. Very common in Spain; parts of Latin America also say \emph{capaz que}.

\begin{cousinweb}
\emph{a lo mejor} is built on \emph{mejor}, inherited from Latin \emph{melior} 'better' (EtymDB); literally the phrase is 'at the best'.
\end{cousinweb}

\begin{grammarlens}[title={What you've built}]
One more word for telling, explaining and describing at B1.
\end{grammarlens}

\subsection*{Your turn}
\begin{itemize}
\raggedright
  \item \emph{Say it:} \emph{a lo mejor}
  \item \emph{Say it:} \emph{a lo mejor}, once more
  \item \emph{Say it:} a sentence of your own with \emph{a lo mejor}
  \item \emph{Recall:} say the Spanish for possibly, then the Spanish for most likely, then say \emph{a lo mejor} again
\end{itemize}

\begin{tcolorbox}[breakable,colback=teal!4,colframe=teal!35!black,title={Before you move on}]
What does a lo mejor mean? (\textquotedblleft{}Maybe, perhaps\textquotedblright{}.) Say it once more.
\end{tcolorbox}
\section[previsible]{previsible — predictable, foreseeable}

\label{lesson:ES-C619-previsible}



\begin{quote}
Before the new one: say the Spanish for most likely, then the Spanish for maybe.
\end{quote}

\subsection*{You'll want to know: previsible}
\textbf{previsible} — \textquotedblleft{}predictable, foreseeable\textquotedblright{}.

\begin{quote}
Era algo previsible.
\end{quote}

\begin{quote}
\emph{It was something you could see coming.}
\end{quote}

The opposite is \emph{imprevisible}.

\begin{cousinweb}
\emph{previsible} is built on \emph{prever}, a learned borrowing from Latin \emph{praevidere} 'to foresee', from \emph{videre} 'to see' (EtymDB).
\end{cousinweb}

\begin{grammarlens}[title={What you've built}]
One more word for telling, explaining and describing at B1.
\end{grammarlens}

\subsection*{Your turn}
\begin{itemize}
\raggedright
  \item \emph{Say it:} \emph{previsible}
  \item \emph{Say it:} \emph{previsible}, once more
  \item \emph{Say it:} a sentence of your own with \emph{previsible}
  \item \emph{Recall:} say the Spanish for most likely, then the Spanish for maybe, then say \emph{previsible} again
\end{itemize}

\begin{tcolorbox}[breakable,colback=teal!4,colframe=teal!35!black,title={Before you move on}]
What does previsible mean? (\textquotedblleft{}Predictable, foreseeable\textquotedblright{}.) Say it once more.
\end{tcolorbox}
\section[la oportunidad]{la oportunidad — the opportunity, the chance}

\label{lesson:ES-C619-oportunidad}



\begin{quote}
Before the new one: say the Spanish for maybe, then the Spanish for predictable.
\end{quote}

\subsection*{You'll want to know: la oportunidad}
\textbf{la oportunidad} — \textquotedblleft{}the opportunity, the chance\textquotedblright{}.

\begin{quote}
Es una buena oportunidad.
\end{quote}

\begin{quote}
\emph{It is a good opportunity.}
\end{quote}

\emph{Aprovechar la oportunidad} — to take the chance.

\begin{cousinweb}
\emph{oportunidad} is a learned borrowing from Latin \emph{opportunitas} 'fitness, a good moment' (EtymDB), built on \emph{opportunus} 'fit, timely', literally 'before the harbour', from \emph{portus} 'harbour' (Lewis \& Short).
\end{cousinweb}

\begin{grammarlens}[title={What you've built}]
One more word for telling, explaining and describing at B1.
\end{grammarlens}

\subsection*{Your turn}
\begin{itemize}
\raggedright
  \item \emph{Say it:} \emph{la oportunidad}
  \item \emph{Say it:} \emph{la oportunidad}, once more
  \item \emph{Say it:} a sentence of your own with \emph{la oportunidad}
  \item \emph{Recall:} say the Spanish for maybe, then the Spanish for predictable, then say \emph{la oportunidad} again
\end{itemize}

\begin{tcolorbox}[breakable,colback=teal!4,colframe=teal!35!black,title={Before you move on}]
What does la oportunidad mean? (\textquotedblleft{}The opportunity, the chance\textquotedblright{}.) Say it once more.
\end{tcolorbox}
\section[la perspectiva]{la perspectiva — the prospect; the perspective}

\label{lesson:ES-C619-perspectiva}



\begin{quote}
Before the new one: say the Spanish for predictable, then the Spanish for the opportunity.
\end{quote}

\subsection*{You'll want to know: la perspectiva}
\textbf{la perspectiva} — \textquotedblleft{}the prospect; the perspective\textquotedblright{}.

\begin{quote}
Las perspectivas de empleo son buenas.
\end{quote}

\begin{quote}
\emph{The job prospects are good.}
\end{quote}

In the plural it often means future chances.

\begin{cousinweb}
\emph{perspectiva} is a learned word usually traced to Medieval Latin \emph{perspectiva} 'the science of optics'; it is built on Latin \emph{perspicere} 'to look through' (Lewis \& Short), from \emph{specere} 'to look'.
\end{cousinweb}

\begin{grammarlens}[title={What you've built}]
One more word for telling, explaining and describing at B1.
\end{grammarlens}

\subsection*{Your turn}
\begin{itemize}
\raggedright
  \item \emph{Say it:} \emph{la perspectiva}
  \item \emph{Say it:} \emph{la perspectiva}, once more
  \item \emph{Say it:} a sentence of your own with \emph{la perspectiva}
  \item \emph{Recall:} say the Spanish for predictable, then the Spanish for the opportunity, then say \emph{la perspectiva} again
\end{itemize}

\begin{tcolorbox}[breakable,colback=teal!4,colframe=teal!35!black,title={Before you move on}]
What does la perspectiva mean? (\textquotedblleft{}The prospect; the perspective\textquotedblright{}.) Say it once more.
\end{tcolorbox}
\section[predecir]{predecir — to predict}

\label{lesson:ES-C619-predecir}



\begin{quote}
Before the new one: say the Spanish for the opportunity, then the Spanish for the prospect.
\end{quote}

\subsection*{You'll want to know: predecir}
\textbf{predecir} — \textquotedblleft{}to predict\textquotedblright{}.

\begin{quote}
Nadie puede predecir el futuro.
\end{quote}

\begin{quote}
\emph{Nobody can predict the future.}
\end{quote}

Like \emph{decir}: \emph{predigo, predije}.

\begin{cousinweb}
\emph{predecir} is a learned borrowing from Latin \emph{praedicere} 'to say beforehand', built from \emph{prae} and \emph{dicere} 'to say', rebuilt on \emph{decir} (EtymDB).
\end{cousinweb}

\begin{grammarlens}[title={What you've built}]
One more word for telling, explaining and describing at B1.
\end{grammarlens}

\subsection*{Your turn}
\begin{itemize}
\raggedright
  \item \emph{Say it:} \emph{predecir}
  \item \emph{Say it:} \emph{predecir}, once more
  \item \emph{Say it:} a sentence of your own with \emph{predecir}
  \item \emph{Recall:} say the Spanish for the opportunity, then the Spanish for the prospect, then say \emph{predecir} again
\end{itemize}

\begin{tcolorbox}[breakable,colback=teal!4,colframe=teal!35!black,title={Before you move on}]
What does predecir mean? (\textquotedblleft{}To predict\textquotedblright{}.) Say it once more.
\end{tcolorbox}
